\begin{proof}[Proof of \cref{obj:thm:matched-curve-frontier}]
Set
\[
\varrho_{n,d}:=\min\left\{1,\frac{d}{n\log(ed)}\right\}.
\]
% lean: CausalSmith.Stat.DiscreteBudgetvalueCurve.curveRate
Choose \(c_\epsilon>0\) from \cref{obj:thm:capacity-active-lower} and \(C_\epsilon^{\mathrm{up}}>0\) from \cref{obj:thm:curve-upper}, and define
\[
C_\epsilon:=C_\epsilon^{\mathrm{up}}+c_\epsilon.
\]
% lean: CausalSmith.Stat.DiscreteBudgetvalueCurve.matched_curve_frontier
The discrete \(L_1\) lower bound underlying \cref{obj:thm:capacity-active-lower} is supplied by \citep[Theorem 3 and the unknown-both-distributions reduction following Theorem 6]{JiaoHanWeissman2018}.

Fix \(n\ge1\), \(d\ge2\), and \(b_0\in[0,1/2]\), and use the independent observed-record laws of \cref{obj:ass:iid-sampling}. Since \(b_0\le1/2\le1-b_0\), evaluation of any admissible curve estimator \(\widehat V\) at \(1/2\) gives a measurable scalar estimator
\[
T_{\widehat V}(O^{(n)}):=\widehat V_{1/2}(O^{(n)}).
\]
% lean: CausalSmith.Stat.DiscreteBudgetvalueCurve.scalar_coordinate_of_curve_estimator
By \cref{obj:thm:capacity-active-lower}, there is a law \(P\in\mathcal M_{d,\epsilon}\) for which
\[
\begin{aligned}
c_\epsilon\varrho_{n,d}
&\le E_P\!\left[\bigl(T_{\widehat V}-V_{1/2}(P)\bigr)^2\right]\\
&\le E_P\!\left[\sup_{b\in[b_0,1-b_0]}
       \bigl|\widehat V_b-V_b(P)\bigr|^2\right]\\
&\le \sup_{P'\in\mathcal M_{d,\epsilon}}
 E_{P'}\!\left[\sup_{b\in[b_0,1-b_0]}
       \bigl|\widehat V_b-V_b(P')\bigr|^2\right].
\end{aligned}
\]
% lean: CausalSmith.Stat.DiscreteBudgetvalueCurve.scalar_coordinate_risk_le_curve_risk
The middle inequality holds samplewise because \(1/2\) belongs to the budget interval. The lower-bound result also supplies a law in the model class when applied to the zero scalar estimator, while \cref{obj:def:jf-estimator} supplies an admissible curve estimator. Thus taking the infimum over curve estimators in \cref{obj:def:curve-risk} gives
\[
c_\epsilon\varrho_{n,d}
\le \mathfrak R_{n,d,\epsilon,b_0}^{\mathrm{curve}}.
\]
% lean: CausalSmith.Stat.DiscreteBudgetvalueCurve.matched_curve_frontier

For the upper bound, insert the estimator of \cref{obj:def:jf-estimator} into that infimum. \Cref{obj:thm:curve-upper} bounds its worst-case curve risk, so
\[
\begin{aligned}
\mathfrak R_{n,d,\epsilon,b_0}^{\mathrm{curve}}
&\le \sup_{P\in\mathcal M_{d,\epsilon}}
 E_P\!\left[\sup_{b\in[b_0,1-b_0]}
 \bigl|\widehat V_b^{\mathrm{JF}}-V_b(P)\bigr|^2\right]\\
&\le C_\epsilon^{\mathrm{up}}\varrho_{n,d}
 \le C_\epsilon\varrho_{n,d}.
\end{aligned}
\]
% lean: CausalSmith.Stat.DiscreteBudgetvalueCurve.matched_curve_frontier
Here \(\varrho_{n,d}\ge0\), and the chosen constants satisfy \(0<c_\epsilon\le C_\epsilon<\infty\).

Now let \(T\) be any measurable scalar estimator. Applying \cref{obj:thm:capacity-active-lower} to the same independent observed-record laws supplies a law \(P\in\mathcal M_{d,\epsilon}\) with \(V_1(P)=1/2\), a binding half-budget policy \(\pi\in\Pi_{1/2}(P)\), and
\[
\sum_{j=1}^d p_j\pi_j=\frac12,\qquad
V_{1/2}(P)=B(P)+\sum_{j=1}^d p_j\pi_j\tau_j,\qquad
E_P\!\left[(T-V_{1/2}(P))^2\right]\ge c_\epsilon\varrho_{n,d}.
\]
% lean: CausalSmith.Stat.DiscreteBudgetvalueCurve.capacity_active_lower
For every occupied cell, the same result gives
\[
p_j>0\quad\Longrightarrow\quad
\frac18\le\tau_j\le\frac38
\quad\Longrightarrow\quad \tau_j>0.
\]
% lean: CausalSmith.Stat.DiscreteBudgetvalueCurve.matched_curve_frontier

Finally, suppose \(b_0=0\) and take any \(P\in\mathcal M_{d,\epsilon}\). Since \(p_j\ge0\) and \(\sum_jp_j=1\), every policy in \([0,1]^d\) is feasible at budget \(1\). For any such policy,
\[
\sum_{j=1}^d p_j\pi_j\tau_j
\le \sum_{j=1}^d p_j\max\{\tau_j,0\},
\]
% lean: CausalSmith.Stat.DiscreteBudgetvalueCurve.full_budget_value
and equality is attained by the policy
\[
\pi_j^*:=\mathbf1\{\tau_j>0\},\qquad j=1,\ldots,d.
\]
% lean: CausalSmith.Stat.DiscreteBudgetvalueCurve.full_budget_value
Consequently, \cref{obj:def:budget-value} gives
\[
V_1(P)=B(P)+\sum_{j=1}^d p_j\max\{\tau_j,0\}.
\]
% lean: CausalSmith.Stat.DiscreteBudgetvalueCurve.full_budget_value
\end{proof}
